Set $A=(0,0),B=(10,0),C=(0,h)$ and $E=(e,0)$. The angle-bisector theorem gives $\dfrac{e}{10-e} = \dfrac{h}{\sqrt{h^2+100}}$, so $h^2 = \dfrac{5 e^2}{5-e}$ and $0<e<5$. The circle with diameter $BD$ has centre $(7,0)$ and radius 3. The distance from its centre to $CE$ gives $h^2(7-e)^2=9(h^2+e^2)$. Substitute the expression for $h^2$ and cancel $e^2>0$ to obtain $5e^2-61e+155=0$. Its only root in $(0,5)$ is $e = \dfrac{61-3\sqrt{69}}{10}$. Therefore $\boxed{B E = \dfrac{39+3\sqrt{69}}{10}}$.
\begin{asy}
import graph;
usepackage("amsmath");
import geometry;
size(7cm);
real e=(61-3sqrt(69))/10,h=sqrt(5e*e/(5-e));pair A=(0,0),B=(10,0),C=(0,h),D=(4,0),E=(e,0);draw(A--B--C--cycle);draw(circle((7,0),3));draw(C--E);label("$A$",A,(-1,-1));label("$B$",B,(1,0));label("$C$",C,(0,1));label("$D$",D,(1,-1));label("$E$",E,(-1,-1));
\end{asy}

